\documentclass[11pt,a4paper]{article}
\usepackage[T1]{fontenc}
\usepackage{lmodern}
\usepackage[margin=25mm]{geometry}
\usepackage{amsmath,amssymb,amsthm}
\usepackage{microtype}
\usepackage{xcolor}
\usepackage{fancyhdr}
\usepackage[colorlinks=true,linkcolor=blue!45!black,pdfauthor={},pdftitle={Finiteness for y cubed minus y equals x to the fourth minus x}]{hyperref}
\definecolor{ink}{HTML}{183D58}
\newtheorem{theorem}{Theorem}
\newtheorem{proposition}{Proposition}
\newcommand{\Z}{\mathbb Z}
\newcommand{\C}{\mathbb C}
\newcommand{\Q}{\mathbb Q}
\newcommand{\PP}{\mathbb P}
\setlength{\parindent}{0pt}
\setlength{\parskip}{6pt}
\pagestyle{fancy}
\fancyhf{}
\fancyhead[L]{\small\color{ink}A Diophantine finiteness correction}
\fancyhead[R]{\small $y^3-y=x^4-x$}
\fancyfoot[C]{\small\thepage}
\setlength{\headheight}{14pt}
\renewcommand{\headrulewidth}{0.3pt}

\begin{document}
\thispagestyle{empty}
\begin{center}
{\Large\bfseries\color{ink}A correction to the proposed infinitude}\\[7pt]
{\LARGE $y^3-y=x^4-x$}\\[8pt]
{\small An unconditional finiteness argument with its theorem dependence explicit}
\end{center}

\textbf{Conclusion and scope.}
The requested assertion is false. The equation has \emph{finitely many}
integer solutions. Below, the geometric calculation is given explicitly,
and the arithmetic conclusion uses Siegel's theorem on integral points.
That theorem is stated but not proved here: this is an unconditional proof,
\emph{not a fully self-contained arithmetic proof}. No proof of the requested
infinitude can exist, and no complete enumeration of the solutions is claimed.
No internet search was used.

\begin{theorem}\label{thm:main}
The set
\[
\{(x,y)\in\Z^2:y^3-y=x^4-x\}
\]
is finite. In fact, its projective algebraic curve is smooth and
geometrically irreducible of genus $3$.
\end{theorem}

\section*{1. Intuition and elementary checks}
The leading terms suggest $y^3\approx x^4$, or $y\approx |x|^{4/3}$.
This makes $x=t^3$, $y=t^4$ a natural first attempt. But substitution gives
\[
(t^4)^3-t^4-\bigl((t^3)^4-t^3\bigr)=t^3(1-t),
\]
which vanishes only for $t=0,1$. Matching the leading powers does not
match the lower terms. The following stronger observation makes this precise.

\begin{proposition}\label{prop:cubes}
If $(x,y)$ is an integer solution and $|x|$ is a perfect cube, then
$x\in\{0,1\}$ and $y\in\{-1,0,1\}$.
\end{proposition}
\begin{proof}
Write $f(s)=s^3-s$. For integers $s\ge1$,
$f(s+1)-f(s)=3s(s+1)>0$.
If $x=\varepsilon t^3$, where $\varepsilon\in\{-1,1\}$ and $t\ge2$,
put $N=x^4-x=t^{12}-\varepsilon t^3$. Then
\begin{align*}
N-f(t^4)&=t^4-\varepsilon t^3>0,\\
f(t^4+1)-N&=3t^8+2t^4+\varepsilon t^3>0.
\end{align*}
Thus $N$ lies strictly between two consecutive values of $f$ on the
positive integers. Since $N>0$ and $f(y)\le0$ for every integer $y\le1$,
no integer $y$ is possible. For $x=-1$, $N=2$ lies between $f(1)=0$
and $f(2)=6$. Finally, $x=0,1$ gives $y(y-1)(y+1)=0$.
\end{proof}

In particular, the six directly verified solutions are
\[
\boxed{(0,-1),\ (0,0),\ (0,1),\ (1,-1),\ (1,0),\ (1,1).}
\]
The proposition does not handle general $x$. To establish global finiteness,
we next determine the geometry of the entire curve.

\newpage
\section*{2. Irreducibility and absence of singularities}
Let the affine curve be $F(x,y)=0$, where
\[
F(x,y)=y^3-y-x^4+x.
\]
We work over $\C$ so that the checks cover all possible components and
singularities, including nonreal ones.

\textbf{Irreducibility.}
As a cubic in $y$ over the field $\C(x)$, a reducible $F$ would have a
root $r(x)=p(x)/q(x)$, with coprime polynomials $p,q\in\C[x]$ and $q\ne0$.
Clearing denominators yields
\[
p^3-pq^2=(x^4-x)q^3.
\]
Hence $q$ divides $p^3$. Coprimality forces $q$ to be constant, so $r$
is a polynomial. A constant $r$ cannot satisfy $r^3-r=x^4-x$; a
polynomial of degree $d\ge1$ would give $3d=4$, also impossible.
Thus $F$ is irreducible over $\C(x)$, and, being monic in $y$, over
$\C[x,y]$ as well.

\textbf{Affine smoothness.}
At an affine singularity, both partial derivatives must vanish:
\[
F_x=1-4x^3=0,\qquad F_y=3y^2-1=0.
\]
Consequently,
\[
x^3=\frac14,\qquad y^2=\frac13.
\]
Substituting $x^4=x/4$ and $y^3=y/3$ into $F=0$ gives
\[
-\frac{2y}{3}+\frac{3x}{4}=0,
\qquad y=\frac98x,
\qquad x^2=\frac{64}{243}.
\]
But the two expressions for $x^6$ would then require
\[
\frac1{16}=\left(\frac{64}{243}\right)^3,
\qquad\text{hence}\qquad 3^{15}=2^{22},
\]
which is impossible. There are no affine singularities.

\textbf{The point at infinity.}
Homogenization gives the projective quartic
\[
H(X,Y,Z)=Y^3Z-YZ^3-X^4+XZ^3=0.
\]
On $Z=0$, this requires $X=0$, leaving exactly the point
$P_\infty=[0:1:0]$. Since
\[
H_Z=Y^3-3YZ^2+3XZ^2,
\qquad H_Z(P_\infty)=1,
\]
this point is smooth too. The homogenization is irreducible: any
projective factorization would yield an affine one unless a factor were
a power of $Z$, and $Z$ does not divide $H$.

We have therefore obtained a smooth, geometrically irreducible projective
curve $C$. Smoothness matters here: a singular quartic can have a much
simpler underlying curve. The checks above exclude that possibility.
The next page computes the genus directly, instead of merely quoting the
genus formula for a smooth plane quartic.

\newpage
\section*{3. A direct genus calculation and the finiteness step}
The coordinate $x$ defines a map $C\to\PP^1(\C)$ of degree $3$, since
for general $x$ the cubic equation has three distinct values of $y$.
Geometric irreducibility makes the associated compact Riemann surface
connected. A point has ramification index $e$ when the map has local form
$t\mapsto t^e$ after local changes of coordinates.

\textbf{Finite ramification.}
Where $F_y\ne0$, the coordinate $x$ is a local parameter, so there is
no ramification. The remaining points have
\[
y=\pm\frac1{\sqrt3},\qquad x^4-x=-\frac{2y}{3}.
\]
For each of these two $y$-values the quartic in $x$ has four distinct
roots: a repeated root would also have $4x^3-1=0$, giving an affine
singularity already excluded. The two sets of $x$-values are disjoint.
At each of these eight points, implicit differentiation with $y$ as local
parameter gives
\[
\frac{dx}{dy}=0,\qquad
\frac{d^2x}{dy^2}=\frac{6y}{4x^3-1}\ne0.
\]
Thus each has ramification index $2$.

\textbf{Ramification at infinity.}
In the chart $Y=1$, put $u=X/Y$, $v=Z/Y$. Near $P_\infty$ the equation is
\[
v-v^3-u^4+uv^3=0,
\qquad v\bigl(1+(u-1)v^2\bigr)=u^4.
\]
Here $u$ is a local parameter because the derivative with respect to $v$
at $(0,0)$ is $1$. The parenthesized factor is nonzero there, so $v$
vanishes to order $4$. Therefore $x=u/v$ has a pole of order $3$,
and the ramification index at infinity is $3$.

\textbf{Counting the genus.}
Triangulate the sphere $\PP^1(\C)$ with all branch values as vertices.
Each open edge and face lifts three times. At a vertex, ramification of
index $e$ reduces the number of preimages by $e-1$ relative to the usual
threefold count. The Euler characteristic is consequently
\[
\chi(C)=3\chi(\PP^1)-\sum_P(e_P-1)
       =6-\bigl(8(2-1)+(3-1)\bigr)=-4.
\]
Since $\chi(C)=2-2g$, the genus is $g=3$.

\textbf{The arithmetic input (Siegel's theorem).}
\emph{If a geometrically irreducible affine algebraic curve over $\Q$
has a smooth projective model of genus at least $1$, then it has only
finitely many integer points in its given affine coordinates.}

Our affine curve is defined over $\Q$, and the preceding computations
give precisely such a model, of genus $3$. Siegel's theorem therefore
proves Theorem~\ref{thm:main}. This step has no unproved conjectural
assumptions, but the theorem itself is a substantial external input.

\textbf{What is and is not established.}
Finiteness is proved; infinitude is disproved. The six displayed points
are verified, but are not proved to be the complete set. No explicit bound
for all solutions, and no self-contained replacement for Siegel's theorem,
is supplied here. These limitations do not change the finiteness conclusion.

\end{document}
